% ============================================================
% 5D MECHANISM BOX TEMPLATE
% For EDC Book 2 - Add throughout chapters
% Date: 2026-01-30
% ============================================================

% ============================================================
% REQUIRED PACKAGE:
% \usepackage{mdframed}
% Or alternative: \usepackage{tcolorbox}
% ============================================================

% ============================================================
% TEMPLATE 1: Using mdframed (simpler)
% ============================================================

\begin{mdframed}[frametitle={\colorbox{white}{\textbf{5D Mechanism: [Title]}}},
                 linewidth=2pt,
                 linecolor=blue!75!black,
                 backgroundcolor=blue!5]

\textbf{Physical Picture in 5D:}

[Describe what's actually happening in 5D geometry]

Example:
% Y-junction with three strings meets at central vertex.
% Strings have tension σ, characteristic length ℓ_p.
% Junction minimizes total energy → Steiner angles (120°).

\textbf{Key Geometric Features:}
\begin{itemize}
\item [Feature 1]: [parameter value with units]
\item [Feature 2]: [parameter value with units]
\item [Feature 3]: [parameter value with units]
\end{itemize}

Example:
% \item String tension: σ = 8.82 MeV/fm² [Dc]
% \item Junction length: ℓ_p = 0.875 fm [Dc]
% \item Angles: θ₁ = θ₂ = θ₃ = 120° (Steiner minimum) [M]

\textbf{Mathematical Framework:}

[Key equations, with brief derivation sketch]

\begin{equation}
[\text{Key equation 1}]
\end{equation}

\begin{equation}
[\text{Key equation 2 if needed}]
\end{equation}

Example:
% E_junction = 3σℓ_p (three strings, equal tension)
% Steiner theorem: ∂E/∂θᵢ = 0 → θᵢ = 120°

\textbf{3D Observable Consequence:}

[What 3D observer sees/measures]

Example:
% Observer measures stable proton with mass m_p.
% Mass arises from junction energy: E_p = m_p c²
% Stability from topological protection (cannot unwind without infinite energy).

\textbf{Validation:}

\begin{tabular}{ll}
Predicted: & [value from 5D] \\
Observed: & [value] [BL] \\
Agreement: & [percentage or "within error bars"] \\
\end{tabular}

Example:
% Predicted: m_p/m_e = 6π⁵ = 1836.118
% Observed: 1836.153 [BL]
% Agreement: 0.002%

\textbf{Epistemic Status}: [\textcolor{blue}{tag}] — [brief explanation]

Example:
% [Dc] — Conditional on Y-junction topology [P] and reduction dictionary

\end{mdframed}

% ============================================================
% TEMPLATE 2: Using tcolorbox (more options)
% ============================================================

\begin{tcolorbox}[colback=blue!5,
                  colframe=blue!75!black,
                  title=\textbf{5D Mechanism: [Title]},
                  fonttitle=\bfseries,
                  boxrule=2pt]

\textbf{Physical Picture in 5D:}

[Same structure as Template 1 above]

\textbf{Key Geometric Features:}
\begin{itemize}
\item [Features as above]
\end{itemize}

\textbf{Mathematical Framework:}
\begin{equation}
[Equations as above]
\end{equation}

\textbf{3D Observable Consequence:}
[As above]

\textbf{Validation:}
[Table as above]

\textbf{Epistemic Status}: [As above]

\end{tcolorbox}

% ============================================================
% EXAMPLE: Thick Brane Necessity
% ============================================================

\begin{mdframed}[frametitle={\colorbox{white}{\textbf{5D Mechanism: Thick Brane Necessity}}},
                 linewidth=2pt,
                 linecolor=blue!75!black,
                 backgroundcolor=blue!5]

\textbf{Physical Picture in 5D:}

The membrane is not infinitely thin but has finite thickness $\delta$.
This creates a ``brane layer'' where 5D $\to$ 3D reduction occurs.
Without thickness, mode overlap integrals diverge and energy cannot be stored temporarily.

\textbf{Key Geometric Features:}
\begin{itemize}
\item Brane thickness: $\delta = \hbar/(2m_p c) = 0.105$ fm [Dc]
\item Energy scale: $E_\delta = \hbar c/\delta \sim m_p c^2$ (proton Compton)
\item Mode localization: Wavefunctions decay as $\psi(y) \sim e^{-|y|/\delta}$
\end{itemize}

\textbf{Mathematical Framework:}

Thin brane ($\delta \to 0$):
\begin{equation}
I_4 = \int dy\, \delta(y) \psi_L(y) \psi_R(y) \quad \text{[DIVERGES]}
\end{equation}

Thick brane ($\delta$ finite):
\begin{equation}
I_4 = \int dy\, \psi_L(y) \psi_R(y) \sim e^{-d_{LR}/\delta} \quad \text{[FINITE]}
\end{equation}

where $d_{LR}$ is chiral separation distance.

\textbf{3D Observable Consequence:}

Fermi coupling suppressed by finite overlap:
\begin{equation}
G_F \sim \frac{g_5^2}{M_W^2} \times I_4
\end{equation}

Finite $\delta$ gives finite $I_4$ $\to$ finite $G_F$.

\textbf{Validation:}

\begin{tabular}{ll}
Framework: & Thick brane with $\delta \sim 0.1$ fm \\
Consequence: & $G_F \sim 10^{-5}$ GeV$^{-2}$ (suppressed) \\
Observed: & $G_F = 1.166 \times 10^{-5}$ GeV$^{-2}$ [BL] \\
Agreement: & Order of magnitude ✓ (precise value pending BVP) \\
\end{tabular}

\textbf{Epistemic Status}: [Dc] — Conditional on $\delta$ scale and mode profile assumptions

\end{mdframed}

% ============================================================
% EXAMPLE: Neutrino Edge Mode
% ============================================================

\begin{tcolorbox}[colback=blue!5,
                  colframe=blue!75!black,
                  title=\textbf{5D Mechanism: Neutrino as Edge Mode},
                  fonttitle=\bfseries,
                  boxrule=2pt]

\textbf{Physical Picture in 5D:}

Neutrino wavefunction $\psi_\nu(y)$ peaks near brane boundary $|y| \sim \delta$
but extends into bulk with exponential tail.
This ``edge mode'' straddles brane-bulk interface, carrying energy between 5D and 3D.

\textbf{Key Geometric Features:}
\begin{itemize}
\item Peak location: $y_{\text{peak}} \sim \pm \delta$ (brane edge)
\item Bulk penetration: $\lambda_{\text{bulk}} \sim$ few $\times \delta$
\item Coupling: Overlaps weakly with brane modes (small $I_\nu$)
\end{itemize}

\textbf{Mathematical Framework:}

Edge mode profile (schematic):
\begin{equation}
\psi_\nu(y) \sim 
\begin{cases}
e^{-(y-\delta)^2/\sigma^2} \times e^{-|y|/\lambda} & y > 0 \\
e^{-(y+\delta)^2/\sigma^2} \times e^{-|y|/\lambda} & y < 0
\end{cases}
\end{equation}

Oscillations from 5D phase evolution:
\begin{equation}
P(\nu_\alpha \to \nu_\beta) \sim \sin^2\left(\frac{\Delta m^2 L}{4E}\right)
\end{equation}

\textbf{3D Observable Consequence:}

From 3D observer perspective:
\begin{itemize}
\item Appears as nearly massless fermion (small $m_\nu$ from bulk tail)
\item Only left-handed couples (5D chirality projection)
\item Oscillates between flavors (5D phase acquired in bulk propagation)
\end{itemize}

\textbf{Validation:}

\begin{tabular}{ll}
Prediction: & Neutrino oscillations (edge mode in 5D) \\
Observed: & Oscillations confirmed (solar, atmospheric, reactor) [BL] \\
Mass ordering: & Normal hierarchy (edge mode energies) \\
Agreement: & Qualitative ✓ (quantitative masses pending) \\
\end{tabular}

\textbf{Epistemic Status}: [P/Dc] — Edge mode is postulated [P]; oscillations are consequence [Dc]

\end{tcolorbox}

% ============================================================
% LIST OF REQUIRED BOXES (for reference)
% ============================================================

% Chapter 1 (3 boxes):
% - Box 1.1: Thick Brane Necessity [EXAMPLE ABOVE]
% - Box 1.2: Junction Relaxation (Neutron)
% - Box 1.3: Energy Ledger + Bulk Exchange

% Chapter 2 (4 boxes):
% - Box 2.1: Particle as Topological Defect
% - Box 2.2: Brane-Dominant Mode
% - Box 2.3: Edge Mode (Neutrino) [EXAMPLE ABOVE]
% - Box 2.4: Composite States (Pions)

% Chapter 3 (2 boxes):
% - Box 3.1: Frozen Projection Operator
% - Box 3.2: Mode Overlap Suppression

% Chapter 4 (3 boxes):
% - Box 4.1: Z₆ → SU(3) Emergence
% - Box 4.2: Discrete Geometry → Continuous Gauge
% - Box 4.3: Three Generation Counting

% Chapters 5-11 (~8 more boxes as needed)

% Total: ~20 boxes throughout book

% ============================================================
% USAGE NOTES:
% ============================================================
% 1. Place boxes AFTER introducing concept, BEFORE detailed math
% 2. Keep physical picture simple (avoid jargon)
% 3. Use consistent color scheme (blue!5 background, blue!75!black frame)
% 4. Always include validation (even if "pending" or "qualitative")
% 5. Tag epistemic status explicitly
% 6. Cross-reference to detailed derivation sections
% ============================================================
